% Non-monotonic comoving-frame radiative transfer in CMFGEN
%
% Compile: module load texlive/2021
%          pdflatex numerical_approach && bibtex numerical_approach && pdflatex numerical_approach && pdflatex numerical_approach

\documentclass[a4paper, 10pt, twocolumn]{article}

\usepackage[round, sort, compress]{natbib}
\usepackage{amsmath, amssymb}
\usepackage[T1]{fontenc}
\usepackage{hyperref}
\usepackage{microtype}
\usepackage[margin=1.8cm]{geometry}

\begin{document}

\title{Non-monotonic comoving-frame radiative transfer in CMFGEN}

\author{S. Brennan\\
\small Max Planck Institute for Astrophysics, Garching, Germany}
\date{}

\maketitle

In the comoving frame, the transfer equation contains a
velocity-coupling term linking adjacent frequencies.  In monotonic
outflows the standard treatment \citep{Hillier1998} absorbs the
coupling into a generalised opacity
$\hat\chi = \chi + \Gamma/\Delta\log\nu$, where
$\Gamma = (V/R)(1+\sigma\mu^2)$ and
$\sigma = (R/V)(\dd V/\dd r) - 1$.  In interacting supernovae the
ejecta contain shocks and compression zones where $\sigma \ll -1$,
making $\hat\chi$ negative and the formal solution unstable.

Following \citet{Knop2009}, we use $|\Gamma|$ in the opacity (keeping
it strictly positive) while retaining the signed $\Gamma$ in the
source term $Q = (\Gamma/\Delta\log\nu)/\hat\chi$.  The upwind
frequency is chosen sign-dependently: previous for expanding flow
($\Gamma > 0$), next for compressing flow ($\Gamma < 0$).  Since the
next-frequency intensity is not yet computed during the sweep, a
Gauss--Seidel scheme supplies it from the previous great iteration.
The moment equations, unstable for $\sigma < -1$ because the Sobolev
resonance has no real root in compression zones, are bypassed;
$J_\nu$ is computed directly from the formal solution.  Each epoch is
treated as a static snapshot with the temperature and electron density
fixed to the hydrodynamic values.  The iteration begins with Lambda
iterations from LTE populations with a maximum correction of $10^5$
\citep{HillierDessart2012}, stable with the direct formal solver, and
automatically transitions to full linearisation when corrections fall
below 50\%.

\bibliographystyle{aasjournal}
\bibliography{references}

\end{document}
